\documentclass[border=10pt]{standalone}
\usepackage[utf8]{inputenc}
\usepackage[T5]{fontenc}
\usepackage[vietnamese]{babel}
\usepackage{tgadventor}
\usepackage{tikz}
\usetikzlibrary{calc, angles, quotes, arrows.meta}

% =====================================================
% Hình minh họa: Tiếp tuyến đường tròn — Bài 3
% O = (0,0), R = 2, A = (-3.5, 0) ngoài đường tròn
% B: tiếp điểm phía trên — OB ⊥ AB
% C: điểm trên đường tròn phía dưới — AC = AB
% angB = acos(-R/OA) = acos(-2/3.5) ≈ 124.85°
% =====================================================

\definecolor{mauchinh} {HTML}{1565C0}
\definecolor{mauphu}   {HTML}{43A047}
\definecolor{maunoibat}{HTML}{E53935}
\definecolor{mauvien}  {HTML}{424242}

% Font qag đậm — dùng cho mọi nhãn
\newcommand{\qagbf}{\fontfamily{qag}\fontseries{b}\selectfont}

\tikzset{
    % Nhãn tên điểm: qag đậm cỡ small
    nhan diem/.style={font=\qagbf\small},
    % Nhãn tên cạnh: qag đậm cỡ scriptsize
    nhan canh/.style={font=\qagbf\scriptsize},
    % Góc vuông
    goc vuong/.style={
        draw=mauvien,
        thick,
        angle radius=0.20cm,
    },
}

\begin{document}
\begin{tikzpicture}[
    declare function={
        % angB: góc OB với trục Ox — điều kiện OB ⊥ AB
        % cos(angB) = -R/OA = -2/3.5
        angB = acos(-2/3.5);
    },
    line join=round,
    line cap=round,
    every node/.style={font=\qagbf\small},
    >=Stealth,
]

% --- Hằng số ---
% R = 2: bán kính đường tròn
\pgfmathsetmacro{\R}{2}

% --- Điểm cơ sở ---
% O: tâm đường tròn
\path (0,0) coordinate (O);
% A: điểm ngoài, trên trục Ox âm
\path (-3.5,0) coordinate (A);

% --- B và C dùng tọa độ cực từ O ---
% B: tiếp điểm phía trên, góc angB
\path (O) ++ ({angB}:\R) coordinate (B);
% C: phía dưới, đối xứng B qua OA, AC = AB
\path (O) ++ ({-angB}:\R) coordinate (C);

% ======= THỨ TỰ VẼ =======

% --- 1. Đường phụ nét đứt ---
% Đoạn OA tham chiếu
\draw[mauvien!50, thin, dashed] (O) -- (A);
% Bán kính OB
\draw[mauchinh!60, thin, dashed] (O) -- (B);
% Bán kính OC
\draw[mauchinh!60, thin, dashed] (O) -- (C);

% --- 2. Đường tròn (O; R) ---
\draw[mauchinh, thick] (O) circle (\R);

% --- 3. Đường chính ---
% Tiếp tuyến AB
\draw[maunoibat, very thick] (A) -- (B);
% Đoạn AC = AB
\draw[mauphu, very thick] (A) -- (C);

% --- 4. Góc vuông tại B: OB ⊥ AB ---
\pic[goc vuong] {right angle = O--B--A};

% --- 5. Vòng lặp: vẽ điểm + nhãn tên đỉnh ---
% \t = tên đỉnh, \col = màu, \g = góc hướng nhãn (tọa độ cực)
\foreach \t/\col/\r/\g in {
    O/mauchinh/1.5pt/-45,
    A/mauvien/1.5pt/180,
    B/maunoibat/1.8pt/125,
    C/mauphu/1.8pt/235
}{
    % Chấm điểm
    \fill[\col] (\t) circle (\r);
    % Nhãn tên đỉnh: shift theo góc cực \g, khoảng 7pt
    \path (\t) node[\col, nhan diem, shift={(\g:7pt)}] {\t};
}

% --- 6. Nhãn tên cạnh — không dùng $...$ ---
% Nhãn R trên bán kính OB
\path (O) -- (B)
    node[midway, nhan canh, mauchinh, shift={(50:7pt)}] {R};
% Nhãn AB
\path (A) -- (B)
    node[midway, nhan canh, maunoibat, shift={(125:7pt)}] {AB};
% Nhãn AC = AB
\path (A) -- (C)
    node[midway, nhan canh, mauphu, shift={(235:7pt)}] {AC = AB};

\end{tikzpicture}
\end{document}
